% ============================================================================
% Fichier  : partie3/P03_C12_cadre_juridique.tex
% Titre    : Le Cadre Juridique Camerounais et Africain
% Auteur   : MINKA MI NGUIDJOÏ Thierry Emmanuel
% Version  : 1.0 -- novembre 2026
% Base   : ACSI C12, repris et développé pour INF3226
% ============================================================================
% Sources et degré de vérification :
%  - Loi n° 2010/012 : texte français consulté (Expertises/.../loi_2010-012_
%    cybersecurite_cybercriminalite.pdf) ; art. 4 (9), 7 (2), 13, 14, 26,
%    29, 32 et 61 cités textuellement.
%  - Règlement COBAC R-2016/04 : art. 1 et 2 repris du cours GCI (relevés
%    sur une base en ligne, à collationner avec la version du SG-COBAC).
%  - AUSCGIE révisé (2014) : missions du commissaire aux comptes et comité
%    d'audit décrits sans citation ; numéros d'articles à vérifier.
%  - Acte uniforme relatif au droit comptable et à l'information financière
%    (2017) : exigences sur la comptabilité informatisée et l'archivage
%    présentées comme A VERIFIER dans le texte.
%  - Loi n° 2018/012 (régime financier) : d'après la présentation du MINFI.
%  - Chambre des comptes, Contrôle Supérieur de l'État, CDBF : rôles
%    généraux ; textes organiques à vérifier.
%  - Loi n° 2024/017 (données personnelles) : d'après le cours GCI, à
%    vérifier au JO.
% ============================================================================

\chapter{Le Cadre Juridique Camerounais et Africain}
\label{chap:droit}
\chaptermark{Cadre juridique}

\epigraph{%
    « L'auditeur constate un écart au critère.\\
    Le juge, seul, qualifie l'infraction. »%
}{Règle de prudence}

\niveauChapitre{Fondamental}{%
    Maîtrise du Chapitre~\ref{chap:normes} (normes professionnelles) et du
    Chapitre~\ref{chap:referentiels} (choix des critères)%
}

\begin{boxobjectifs}
\begin{itemize}
    \item Identifier les obligations d'audit et de contrôle du droit
          camerounais et africain~: loi n°~2010/012, règlement COBAC
          R-2016/04, droit OHADA des sociétés.
    \item Situer les exigences du droit comptable OHADA sur la comptabilité
          informatisée et l'archivage, en distinguant ce qui est vérifié de
          ce qui reste à vérifier.
    \item Situer le contrôle des finances publiques~: régime financier de
          l'État, Chambre des comptes, Contrôle Supérieur de l'État.
    \item Analyser la responsabilité de l'auditeur, son secret
          professionnel et la limite entre constat d'audit et qualification
          juridique.
\end{itemize}
\end{boxobjectifs}

\bigskip

% ============================================================================
\section*{Question directrice : que doit l'auditeur à la loi~?}
% ============================================================================

Trois scènes, la même semaine. Chez LiZbiZ, le directeur général demande à
AUDIX CONSEIL de « ne pas ébruiter » les ajustements de stock suspects
du Chapitre~\ref{chap:fraude}, et de remettre le rapport au seul conseil.
À la BANQUE SANAGA, l'ANTIC annonce son audit de sécurité annuel, et le RSSI
se demande ce qu'il doit remettre. À l'ANPP, la Chambre des comptes relève
des paiements sans pièce justificative, et M.~NJOYA se demande s'il doit
écrire « paiements irréguliers » ou « paiements illégaux ».

\medskip
\textbf{Plan du chapitre.}
\begin{enumerate}
    \item Le panorama des textes (section~\ref{sec:c15-panorama}).
    \item La loi sur la cybersécurité et la loi sur les communications électroniques (section~\ref{sec:c15-2010}).
    \item La réglementation bancaire de la CEMAC (section~\ref{sec:c15-cobac}).
    \item Le droit OHADA (section~\ref{sec:c15-ohada}).
    \item Le contrôle des finances publiques (section~\ref{sec:c15-public}).
    \item La responsabilité et le secret de l'auditeur (section~\ref{sec:c15-responsabilite}).
    \item Transformer un texte en critères d'audit (section~\ref{sec:c15-criteres}).
\end{enumerate}

\begin{boxalert}
\textbf{Règle du chapitre.} Un texte de loi n'est cité entre guillemets
que s'il a été lu dans sa version officielle. Les autres dispositions sont
présentées sans citation, avec la mention «~à vérifier~». Un auditeur
applique la même règle dans ses rapports.
\end{boxalert}

\separateur

% ============================================================================
\section{Panorama des Textes}
\label{sec:c15-panorama}
% ============================================================================

\begin{tableaucours}{L{3.0cm}L{2.8cm}L{4.0cm}L{2.2cm}}{%
    \tableauHeader{Texte} &
    \tableauHeader{Assujettis} &
    \tableauHeader{Ce qu'il impose à l'audit} &
    \tableauHeader{Vérification}
}
Loi n°~2010/012 (cybersécurité) & Opérateurs, exploitants de SI, autorités
de certification & Audit de sécurité obligatoire~; secret des
auditeurs~; sanction de l'obstruction & Texte officiel lu \\
Règlement COBAC R-2016/04 & Établissements de crédit de la CEMAC &
Dispositif de contrôle interne, dont le risque informatique & Art.~1 et~2
relevés \\
AUSCGIE révisé (2014) & Sociétés commerciales & Commissaire aux comptes~;
comité d'audit des sociétés cotées & Présenté, à vérifier \\
Acte uniforme sur le droit comptable (2017) & Entités qui tiennent une
comptabilité OHADA & Organisation comptable, comptabilité informatisée,
conservation & Présenté, à vérifier \\
Loi n°~2010/013 (communications électroniques) & Opérateurs et
fournisseurs de services de communications électroniques & Obligations des
opérateurs, confidentialité des communications & Présenté, à vérifier \\
Loi n°~2018/012 (régime financier) & État et entités publiques & Contrôle
interne, responsables de programme & D'après le MINFI \\
Loi n°~2024/017 (données personnelles) & Responsables de traitement &
Traitement des données lors de l'audit & À vérifier au JO \\
\end{tableaucours}
\vspace{-0.8em}
\begin{center}{\small\itshape Tableau~15.1. Les textes qui encadrent l'audit
des SI au Cameroun, et leur degré de vérification dans ce cours}\end{center}

% ============================================================================
\section{La Loi sur la Cybersécurité}
\label{sec:c15-2010}
% ============================================================================

\subsection{L'audit de sécurité obligatoire}

La loi n°~2010/012 relative à la cybersécurité et à la cybercriminalité
définit l'audit de sécurité et le rend obligatoire pour certaines
organisations~\cite{loi2010012}.

\begin{boxjuris}
\textbf{\loicam{n°~2010/012}, article~4, point~9}~: \textit{« Audit de
sécurité~: examen méthodique des composantes et des acteurs de la sécurité,
de la politique, des mesures, des solutions, des procédures et des moyens
mis en oeuvre par une organisation, pour sécuriser son environnement,
effectuer des contrôles de conformité, des contrôles d'évaluation de
l'adéquation des moyens (organisationnels, techniques, humains,
financiers) investis au regard des risques encourus, d'optimisation, de
rationalité et de performance. »}

\medskip
\textbf{Article~13, alinéa~1}~: \textit{« Sont soumis à un audit de sécurité
obligatoire, les réseaux de communications électroniques et les systèmes
d'information, des opérateurs, des autorités de certification et des
fournisseurs de services de communications électroniques. »}

\medskip
\textbf{Article~32}~: \textit{« (1) Les réseaux de communications
électroniques et les systèmes d'information sont soumis à un régime
d'audit de sécurité obligatoire et périodique de leurs systèmes de
sécurité par l'Agence. (2) L'audit de sécurité et les mesures d'impact de
gravité sont effectués chaque année ou lorsque les circonstances
l'exigent. (3) Les rapports d'audit sont confidentiels et adressés au
Ministre chargé des Télécommunications. »}

\medskip
\textbf{Analyse.} La définition légale couvre à la fois la conformité et
l'adéquation des moyens aux risques~: elle rejoint l'audit de conformité
\emph{et} l'audit d'efficacité du Chapitre~\ref{chap:audit-auditologie}.
L'audit est annuel ou déclenché par les circonstances, et son rapport est
confidentiel. Les conditions et modalités sont renvoyées à des textes
réglementaires (art.~13, al.~2), que l'auditeur doit consulter.
\end{boxjuris}

L'ANTIC, régulateur de la sécurité électronique, a notamment pour mission,
selon l'article~7, alinéa~2, \textit{« de s'assurer de la régularité, de
l'effectivité des audits de sécurité des systèmes d'information suivant
les normes en la matière, des organismes publics et des autorités de
certification »}. La loi renvoie donc aux normes~: un audit de sécurité
légal se conduit selon les normes du Chapitre~\ref{chap:normes} et les
référentiels du Chapitre~\ref{chap:referentiels}.

\subsection{Les obligations qui fournissent des critères}

Plusieurs articles imposent aux \emph{exploitants des systèmes
d'information} des obligations qui deviennent, pour l'auditeur, des
critères opposables.

\begin{tableaucours}{L{1.6cm}L{6.4cm}L{4.2cm}}{%
    \tableauHeader{Art.} &
    \tableauHeader{Obligation (texte officiel, résumé)} &
    \tableauHeader{Critère d'audit}
}
26 (1) & Prendre les mesures techniques et administratives de sécurité~;
se doter de systèmes normalisés pour identifier, évaluer, traiter et gérer
continûment les risques & Un processus de gestion des risques de sécurité
existe et fonctionne \\
26 (2) et (3) & Mettre en place des mécanismes contre les atteintes à la
disponibilité, l'intégrité, l'authentification, la non-répudiation, la
confidentialité et la sécurité physique, approuvés et visés par l'Agence &
Les mécanismes existent~; le visa de l'Agence est au dossier \\
29 (1) & Conserver les données de connexion et de trafic pendant dix ans &
La politique de conservation couvre dix ans \\
29 (2) & Installer des mécanismes de surveillance du contrôle d'accès aux
données & Les accès aux données sont journalisés et surveillés \\
30 (1) & Évaluer et réviser les systèmes de sécurité selon l'évolution des
technologies & Une revue périodique est tracée \\
\end{tableaucours}
\vspace{-0.8em}
\begin{center}{\small\itshape Tableau~15.2. Obligations de la loi
n°~2010/012 et critères d'audit correspondants}\end{center}

\begin{boxlecture}{Article 29, LiZbiZ}
AUDIX CONSEIL ajoute au domaine~D quatre contrôles tirés de la loi~: visa de
l'Agence sur les mécanismes de sécurité (art.~26, al.~3), conservation des
journaux de connexion pendant dix ans (art.~29, al.~1), surveillance des
accès aux données (art.~29, al.~2), revue périodique des systèmes de
sécurité (art.~30). Cotation~: 1~\OUI{} (surveillance des accès), 2~\NON{}
(les journaux sont conservés douze mois~; aucune revue n'est tracée),
1~\NEK{} (le DSI affirme avoir obtenu le visa, sans pièce).

\medskip
\textbf{Constatation.} $\mu = 0{,}25$~; $U = 0{,}25$~; $H = 1{,}50$~bit.

\medskip
\textbf{Analyse} (confiance moyenne). Non-conformité opaque sur un petit
ensemble. Les deux \NON{} sont des constats légaux~: le rapport cite
l'article, et la recommandation sur la conservation pèse sur le budget de
stockage des quatre centres de données. Le \NEK{} se lève par une simple
demande de la pièce à l'Agence ou au DSI.

\medskip
\hfill\lecture{0{,}25}{0{,}25}{1{,}50}
\end{boxlecture}

\begin{boxmetier}{RSSI}
\textbf{Préparer l'audit de sécurité légal}
\begin{enumerate}
    \item Établir si l'organisation est assujettie~: opérateur, autorité de
          certification, fournisseur de services, exploitant de SI.
    \item Rassembler les pièces des articles~26, 29 et~30~: analyse des
          risques, mécanismes de sécurité et visa de l'Agence, politique de
          conservation, revues de sécurité.
    \item Consulter les textes réglementaires d'application (art.~13,
          al.~2, et art.~32, al.~4) en vigueur à la date de l'audit.
    \item Désigner l'interlocuteur des auditeurs et répondre aux demandes
          dans les délais~: l'obstruction est une infraction (art.~61,
          al.~3).
    \item Suivre les recommandations jusqu'à l'audit suivant, qui est
          annuel.
\end{enumerate}
\end{boxmetier}

% ============================================================================
\section{La Réglementation Bancaire de la CEMAC}
\label{sec:c15-cobac}
% ============================================================================

Les établissements de crédit de la CEMAC sont soumis au règlement COBAC
R-2016/04 relatif au contrôle interne~\cite{cobac2016r04}. Le cours GCI en a
relevé les deux premiers articles.

\begin{boxjuris}
\textbf{Règlement COBAC R-2016/04, articles~1 et~2} (texte relevé sur une
base réglementaire en ligne, à collationner avec la version du
SG-COBAC)~:

\medskip
Article~1~: \textit{« Le présent règlement fixe les exigences minimales
relatives au système de contrôle interne que les établissements assujettis
doivent mettre en œuvre. »}

\medskip
Article~2, risque opérationnel~: \textit{« risque de pertes résultant de
l'inadéquation ou de la défaillance des processus internes, des personnes
et des systèmes, ou d'évènements externes. »}

\medskip
\textbf{Analyse.} Une défaillance informatique est un risque opérationnel
au sens du règlement. L'évaluation du contrôle interne de la BANQUE SANAGA
(Chapitre~\ref{chap:referentiels}) répond donc à une exigence
réglementaire, et la panne de mars~2026 en est un incident à documenter.
\end{boxjuris}

Pour l'auditeur, le règlement fournit le cadre~; le COSO fournit la
méthode d'évaluation, et COBIT ou l'ISO/IEC~27002 les critères
informatiques. Le rapport à la COBAC dit, pour chaque déficience, quel
risque opérationnel elle laisse sans maîtrise.

% ============================================================================
\section{Le Droit OHADA}
\label{sec:c15-ohada}
% ============================================================================

\subsection{Le commissaire aux comptes et le comité d'audit}

L'Acte uniforme relatif au droit des sociétés commerciales et du
groupement d'intérêt économique (\sigle{AUSCGIE}), révisé le 30~janvier
2014, organise la mission du commissaire aux
comptes~\cite{ohada2014auscgie}. Selon la doctrine, cette mission
comprend notamment~:
\begin{itemize}
    \item la certification de la régularité, de la sincérité et de l'image
          fidèle des états financiers~;
    \item la vérification permanente des valeurs et des documents
          comptables, sans immixtion dans la gestion~;
    \item le signalement à l'assemblée des irrégularités et des
          inexactitudes relevées~;
    \item la révélation au ministère public des faits délictueux dont il a
          eu connaissance dans l'exercice de sa mission.
\end{itemize}
L'Acte révisé rend en outre obligatoire, dans les sociétés cotées, un
comité d'audit chargé notamment du suivi de l'efficacité des systèmes de
contrôle interne et de gestion des risques. Les numéros d'articles sont à
vérifier dans le texte avant toute citation.

\begin{boxinfo}
\textbf{Le commissaire aux comptes.}
\textbf{Les indices de fraude de LiZbiZ}
\begin{enumerate}
    \item Apprécier l'effet des ajustements de stock suspects sur les états
          financiers~: stock surévalué, résultat faussé.
    \item Signaler les irrégularités relevées à l'organe compétent, puis à
          l'assemblée.
    \item Si les faits paraissent délictueux, les révéler au ministère
          public~: la demande du directeur général de « ne pas ébruiter »
          ne le délie pas de cette obligation.
    \item Documenter dans son dossier les éléments qui fondent cette
          appréciation, dont les travaux d'AUDIX CONSEIL.
\end{enumerate}
\end{boxinfo}

AUDIX CONSEIL, auditeur contractuel, n'a pas cette obligation de
révélation~: il remet son rapport à son commanditaire, selon la lettre de
mission. Mais il ne peut pas retirer de son rapport un constat qu'il a
établi (Chapitre~\ref{chap:mission}), et le rapport sera lu par le
commissaire aux comptes.

\subsection{La comptabilité informatisée et l'archivage}

L'Acte uniforme relatif au droit comptable et à l'information financière,
adopté en 2017, fixe l'organisation comptable des entités de l'espace
OHADA. Les exigences suivantes intéressent directement l'auditeur des SI~;
elles sont présentées d'après le plan du cours et \emph{restent à vérifier
dans le texte} avant toute citation~:
\begin{itemize}
    \item une organisation comptable documentée, qui décrit aussi les
          traitements informatisés~;
    \item des écritures qui, une fois validées, ne peuvent plus être
          modifiées ni supprimées, les corrections passant par de nouvelles
          écritures~;
    \item la possibilité de remonter de chaque montant des états financiers
          à la pièce justificative, et inversement (piste d'audit)~;
    \item la conservation des documents comptables et des pièces
          justificatives pendant une durée minimale fixée par le texte.
\end{itemize}

\begin{boiteRemarque}{Pourquoi signaler « à vérifier »~?}
Parce qu'un critère juridique mal cité rend le constat contestable
(axiome~6). L'auditeur qui écrit « l'Acte uniforme impose
l'irréversibilité des écritures, article~X » doit avoir lu l'article~X.
S'il ne l'a pas lu, il écrit son constat sur un critère qu'il maîtrise,
par exemple la procédure comptable de l'entité, et signale la question
juridique à vérifier.
\end{boiteRemarque}

% ============================================================================
\section{Le Contrôle des Finances Publiques}
\label{sec:c15-public}
% ============================================================================

\subsection{Le régime financier de l'État}

La loi n°~2018/012 du 11~juillet~2018 portant régime financier de l'État et
des autres entités publiques organise le budget par programmes. Selon la
présentation du ministère des Finances, elle institue des responsables de
programme, qui veillent au respect des dispositifs de contrôle interne et
de contrôle de gestion~\cite{minfi2023regimefinancier}. La dépense publique
reste soumise à la distinction classique entre l'ordonnateur, qui engage et
ordonne la dépense, et le comptable public, qui la paie après contrôle~;
dans un système de paiement informatisé, cette séparation se vérifie dans
la matrice des droits, comme au Chapitre~\ref{chap:controle-interne}.

\subsection{Les organes de contrôle}

\begin{tableaucours}{L{3.4cm}L{3.6cm}L{5.2cm}}{%
    \tableauHeader{Organe} &
    \tableauHeader{Nature du contrôle} &
    \tableauHeader{Ce qui intéresse l'auditeur des SI}
}
Chambre des comptes de la Cour suprême & Juridictionnel, sur les comptes
publics & Audits de conformité et de performance~; normes ISSAI
(Chapitre~\ref{chap:normes}) \\
Contrôle Supérieur de l'État & Administratif, rattaché à la Présidence de
la République & Missions de vérification des services et organismes
publics \\
Conseil de discipline budgétaire et financière & Sanction des fautes de
gestion des agents publics & Suites possibles des constats d'audit \\
Contrôleurs financiers et inspections ministérielles & Contrôle interne de
la dépense & Premier niveau de contrôle des paiements \\
\end{tableaucours}
\vspace{-0.8em}
\begin{center}{\small\itshape Tableau~15.3. Les organes de contrôle des
finances publiques (rôles généraux~; textes organiques à vérifier)}\end{center}

\begin{boxlecture}{Paiements de l'ANPP}
M.~NJOYA examine un échantillon aléatoire de 40~paiements de l'ANPP au
regard de quatre critères~: pièce justificative au dossier, bénéficiaire
en activité, ordonnateur distinct du comptable, visa du contrôle
financier. Un paiement est \OUI{} s'il satisfait les quatre.

\medskip
Résultat~: 29~\OUI{}~; 6~\NON{}, dont 4~paiements à des agents partis
dont les comptes restaient actifs~; 5~\NEK{}, dont la pièce n'est pas au
dossier communiqué, l'agence affirmant qu'elle se trouve aux archives.

\medskip
\textbf{Constatation.} $\mu \approx 0{,}73$~; $U \approx 0{,}13$~;
$H \approx 1{,}12$~bit.

\medskip
\textbf{Analyse} (confiance moyenne). Les cinq \NEK{} ne sont pas des
\NON{}~: tant que l'agence peut produire les pièces, on ne peut pas
conclure à leur absence. Le magistrat fixe un délai écrit~; une pièce non
produite à l'échéance, avec confirmation écrite de son absence, devient
\NON{}. Les quatre paiements à des agents partis relient le contrôle des
paiements au contrôle d'accès~: c'est le même défaut, vu de deux côtés. Le
ACSI-LAB-06 reprend ce cas.

\medskip
\hfill\lecture{0{,}73}{0{,}13}{1{,}12}
\end{boxlecture}

% ============================================================================
\section{Transformer un Texte en Critères}
\label{sec:c15-criteres}
% ============================================================================

Un texte de loi ne se cote pas tel quel~: il contient souvent plusieurs
obligations dans un même article, et des termes qui demandent à être
précisés (« mesures appropriées », « dans les meilleurs délais »).
L'auditeur le transforme en critères atomiques
(Chapitre~\ref{chap:decomposition}).

\begin{boxprocedure}
\textbf{Du texte au critère, en quatre étapes}
\begin{enumerate}
    \item \textbf{Lire le texte officiel}, et noter sa date et sa source
          (Journal officiel, recueil de la COBAC, site de l'OHADA).
    \item \textbf{Isoler chaque obligation}~: un sujet (qui est tenu~?), un
          verbe (que doit-il faire~?), un objet, éventuellement un délai.
    \item \textbf{Préciser les termes vagues} par un texte d'application,
          une norme ou une décision de l'organisation, et noter la source de
          la précision.
    \item \textbf{Écrire le critère et la question}, un par obligation, et
          les rattacher à la fiche d'assertion du contrôle.
\end{enumerate}
\end{boxprocedure}

\begin{boxscenario}{Opération MEFOU. Un article, trois critères}
\textbf{Texte} (régime financier, présenté d'après le cours~; à vérifier
dans le texte officiel)~: la dépense publique passe par l'engagement, la
liquidation, l'ordonnancement et le paiement, et l'ordonnateur et le
comptable sont des personnes distinctes.

\medskip
\textbf{Critères tirés par M.~NJOYA.} (1) Chaque paiement est précédé d'un
ordonnancement daté~; (2) l'ordonnateur et le comptable d'un même paiement
sont deux personnes différentes~; (3) chaque paiement porte le visa du
contrôle financier antérieur à son exécution.

\medskip
\textbf{Constatation} (40~paiements, Chapitre~\ref{chap:audit-conformite}).
Un paiement a le même agent pour ordonnateur et pour comptable~: le
critère~(2) est \NON{} pour ce paiement, quel que soit le résultat des deux
autres.
\end{boxscenario}

% ============================================================================
\section{Responsabilité et Secret de l'Auditeur}
\label{sec:c15-responsabilite}
% ============================================================================

\subsection{Une obligation de moyens}

L'auditeur contractuel s'engage, dans sa lettre de mission, à conduire ses
travaux selon les normes qu'il déclare. Il n'est pas garant de l'absence
d'anomalie~: son engagement est, en règle générale, une obligation de
moyens. Sa responsabilité peut être recherchée s'il n'a pas appliqué ces
normes~: d'où l'importance des papiers de travail
(Chapitre~\ref{chap:preuve}), qui montrent ce qu'il a fait, et des \NEK{}
documentés, qui montrent ce qu'il n'a pas pu établir.

\subsection{Le secret professionnel}

\begin{boxjuris}
\textbf{\loicam{n°~2010/012}, article~14}~: \textit{« Le personnel de
l'Agence et les experts commis en vue d'accomplir des opérations d'audit
sont astreints au secret professionnel. »}

\medskip
\textbf{Article~61, alinéa~1}~: \textit{« Sont punis d'un emprisonnement de
trois (03) mois à trois (03) ans et d'une amende de 20.000 (vingt mille) à
100.000 (cent mille) F CFA, les personnels de l'Agence et les experts des
personnes morales chargés des audits qui révèlent sans autorisation, des
informations confidentielles dont ils ont eu connaissance à l'occasion
d'un audit de sécurité. »}

\medskip
\textbf{Article~61, alinéa~3}~: est puni d'un emprisonnement d'un à cinq
ans et d'une amende, ou de l'une de ces deux peines, celui qui fait
obstacle aux audits de sécurité ou refuse de fournir les informations ou
documents y afférents.

\medskip
\textbf{Analyse.} La loi protège l'audit dans les deux sens~: elle punit
l'auditeur qui divulgue, et l'audité qui fait obstruction. Pour un audit
de sécurité légal, un refus de remettre un document n'est donc pas
seulement un \NEK{}~: c'est aussi un fait que l'auditeur doit consigner.
\end{boxjuris}

Le secret couvre les failles constatées, les données consultées et le
rapport lui-même. Il ne fait pas obstacle aux obligations de révélation
que la loi impose à certains auditeurs, comme le commissaire aux comptes.
Lorsque l'auditeur extrait des données personnelles (fichiers du
personnel, bénéficiaires de paiements), il ne prélève que ce que le test
exige, les protège et les détruit à la fin de la conservation du dossier,
dans le respect de la loi n°~2024/017 sur la protection des données
personnelles (texte à vérifier au Journal officiel).

\subsection{Constater n'est pas qualifier}

\begin{boxalert}
\textbf{Violation typique.} Écrire dans un rapport d'audit « M.~FOUDA a
commis une fraude » ou « ces paiements sont illégaux ».

\smallskip
\textbf{Règle~:} l'auditeur constate un écart à un critère, preuve à
l'appui~: « 148~millions FCFA d'ajustements de stock, dont 71~\% en fin
de mois, enregistrés sous le compte de M.~FOUDA ». La qualification pénale
(fraude, détournement) appartient au juge. M.~NJOYA écrit « paiements non
conformes au régime financier », et la juridiction compétente en tire les
conséquences.
\end{boxalert}

\begin{boxmetier}{Auditeur SI}
\textbf{Protéger la mission et les personnes}
\begin{enumerate}
    \item Vérifier si la mission relève d'un régime légal (audit de
          sécurité de la loi n°~2010/012, contrôle de la COBAC, contrôle
          public) et en tirer les obligations de forme et de destinataire.
    \item Inscrire dans la lettre de mission la confidentialité, l'accès
          aux documents et le sort des données extraites.
    \item Rédiger les constats en termes d'écart au critère, sans
          qualification juridique.
    \item Signaler au commanditaire, par écrit, tout fait qui pourrait
          appeler une révélation par un tiers tenu de la faire.
    \item Consigner tout refus de communication, avec sa date et son
          auteur.
\end{enumerate}
\end{boxmetier}

% ============================================================================

\separateur

\begin{boxresume}
\textbf{Ce qu'il faut retenir de ce chapitre~:}
\begin{itemize}
    \item \textbf{La loi n°~2010/012} définit l'audit de sécurité, le rend
          obligatoire et annuel pour les assujettis (art.~13 et~32), et
          fournit des critères opposables (art.~26, 29 et~30).
    \item \textbf{Elle protège l'audit}~: secret des auditeurs (art.~14),
          sanctions de la divulgation et de l'obstruction (art.~61).
    \item \textbf{Le règlement COBAC R-2016/04} fait de la défaillance
          informatique un risque opérationnel, couvert par le contrôle
          interne exigé des banques.
    \item \textbf{Le droit OHADA} organise le commissaire aux comptes, tenu
          de révéler les faits délictueux, et le comité d'audit~; ses
          exigences sur la comptabilité informatisée sont à vérifier dans
          le texte.
    \item \textbf{Le contrôle public} relève de la Chambre des comptes, du
          Contrôle Supérieur de l'État et du CDBF~; un \NEK{} sur une pièce
          justificative se lève dans un délai écrit.
    \item \textbf{L'auditeur constate}, le juge qualifie~; on ne cite un
          texte que si on l'a lu.
    \item \textbf{Un texte se transforme en critères atomiques}~: lire la version officielle, isoler chaque obligation, préciser les termes vagues, citer la source.
\end{itemize}
\end{boxresume}

\bigskip

\begin{boxexo}[title={\ding{45}\quad Exercice~15.1. Assujetti ou non~? \niveaubase}]
Indiquez, pour chaque organisation, si elle relève de l'audit de sécurité
obligatoire de la loi n°~2010/012, et des autres textes du chapitre.
\begin{enumerate}[(a)]
    \item Un opérateur de téléphonie mobile.
    \item La BANQUE SANAGA.
    \item LiZbiZ, société de distribution qui exploite quatre centres de
          données.
    \item L'ANPP.
\end{enumerate}
\end{boxexo}

\begin{boxexo}[title={\ding{45}\quad Exercice~15.2. Rédiger un constat légal \niveaumoyen}]
Reprenez la lecture « Article 29, LiZbiZ ».
\begin{enumerate}[(a)]
    \item Rédigez l'observation relative à la conservation des journaux, en
          sept éléments (Chapitre~\ref{chap:conclure}).
    \item Quelle pièce lèverait le \NEK{} sur le visa de l'Agence~?
    \item Si le DSI refuse de la produire, quelle conséquence juridique
          cela a-t-il dans le cadre d'un audit de sécurité légal~?
    \item Proposez une recommandation réaliste sur la conservation, en
          tenant compte de la loi sur les données personnelles.
\end{enumerate}
\end{boxexo}

\begin{boxexo}[title={\ding{45}\quad Exercice~15.3. Taire ou révéler~? \niveauexpert}]
Le directeur général de LiZbiZ demande à AUDIX CONSEIL de retirer du rapport
les constats sur les ajustements de stock de M.~FOUDA, et de les traiter
oralement.
\begin{enumerate}[(a)]
    \item AUDIX CONSEIL peut-il accepter~? Fondez votre réponse sur les
          normes (Chapitre~\ref{chap:normes}) et sur l'axiome~6.
    \item Quelle est la position du commissaire aux comptes de LiZbiZ s'il
          a connaissance de ces faits~?
    \item Rédigez le constat sans qualification juridique.
    \item Le secret professionnel d'AUDIX CONSEIL l'empêche-t-il
          d'informer le commissaire aux comptes~? Discutez.
\end{enumerate}
\end{boxexo}

\bigskip

\begin{boxexo}[title={\ding{45}\quad Exercice~15.4. Quel texte~? \niveaubase}]
Pour chacune des situations, dites quel texte du Tableau~15.1 s'applique~:
(a) un opérateur de téléphonie conserve les données de connexion~; (b) la
BANQUE SANAGA n'a aucun contrôle sur les changements~; (c) LiZbiZ envoie
les données de ses clients à un prestataire étranger~; (d) un paiement de
l'ANPP n'a pas de visa.
\end{boxexo}

\begin{boxexo}[title={\ding{45}\quad Exercice~15.5. Du texte au critère \niveaumoyen}]
Appliquez la procédure en quatre étapes à l'obligation, présentée par le
cours, de sécuriser les systèmes d'information (loi n°~2010/012). Écrivez
au moins trois critères atomiques et leur question.
\end{boxexo}

\begin{boxexo}[title={\ding{45}\quad Exercice~15.6. Citer « à vérifier » \niveaumoyen}]
Un auditeur junior cite dans son rapport un article de l'AUSCGIE trouvé sur
un forum, entre guillemets, avec un numéro d'article. Que faites-vous~?
Rédigez la phrase corrigée.
\end{boxexo}

\begin{boxexo}[title={\ding{45}\quad Exercice~15.7. Irrégulier ou illégal~? \niveaumoyen}]
M.~NJOYA hésite entre « paiements irréguliers » et « paiements illégaux ».
Que conseillez-vous, et pourquoi~? Rédigez la constatation sur les quatre
agents payés après leur départ.
\end{boxexo}

\begin{boxexo}[title={\ding{45}\quad Exercice~15.8. Le secret et la réquisition \niveaumoyen}]
Un officier de police judiciaire demande à AUDIX CONSEIL ses papiers de
travail sur LiZbiZ. Quelles questions l'associé doit-il se poser avant de
répondre~? Quelle part de la réponse relève du conseil d'un avocat~?
\end{boxexo}

\begin{boxexo}[title={\ding{45}\quad Exercice~15.9. Un audit ANTIC \niveauexpert}]
L'ANTIC annonce son audit annuel de la BANQUE SANAGA. Le RSSI vous demande
de préparer le dossier. Listez les pièces à remettre, celles que la
banque peut refuser et sur quel fondement, et la manière de traiter un
constat de l'ANTIC qui contredit le rapport de l'audit interne.
\end{boxexo}

\begin{boxexo}[title={\ding{45}\quad Exercice~15.10. Données personnelles et papiers de travail \niveauexpert}]
Les extractions de l'audit de LiZbiZ contiennent les noms, numéros de
téléphone et adresses de 12~000~clients. Proposez des règles de traitement
de ces données par AUDIX CONSEIL (minimisation, durée, sécurité,
destruction), en indiquant ce qui relève de la loi n°~2024/017 et ce qui
relève de la déontologie.
\end{boxexo}

\bigskip

\begin{boxinfo}
\textbf{Points de contrôle de l'élève.} Je sais~: citer les textes qui
encadrent l'audit des SI au Cameroun et leur degré de vérification~;
expliquer l'audit de sécurité obligatoire~; dire ce que le COBAC impose~;
distinguer constater et qualifier~; transformer un article en critères
atomiques~; appliquer la règle « citer seulement ce qu'on a lu ».
\end{boxinfo}

\begin{boxinfo}
\textbf{Pour aller plus loin.}
\begin{itemize}
    \item Loi n°~2010/012 relative à la cybersécurité et à la
          cybercriminalité~\cite{loi2010012}.
    \item COBAC, règlement R-2016/04~\cite{cobac2016r04}.
    \item OHADA, AUSCGIE révisé~\cite{ohada2014auscgie}.
    \item MINFI, régime financier de l'État~\cite{minfi2023regimefinancier}.
\end{itemize}
\end{boxinfo}

\begin{boxinfo}
\textbf{Transition.} La loi fixe ce que l'auditeur doit faire et ce qu'il
doit taire~; la déontologie fixe la manière dont il doit se conduire, y
compris là où la loi se tait. Le Chapitre~\ref{chap:deontologie} présente
les codes de la profession, la responsabilité de l'auditeur, ses
certifications et le marché de l'audit.
\end{boxinfo}
